
où \(A\) est l'%
anneau \(K[X_1,\ldots,X_n]\), filtré par le poids par rapport aux \(X_i\), \(X_i\) étant de poids \(i\), et pour l'élément \((n,1+X_1+\ldots+X_n)\).

On plonge \(A\) dans \(A'=K[T_1,\ldots,T_n]\) par les fonctions symétriques, et on obtient ainsi un homomorphisme injectif : \(\Ch A\longrightarrow\Ch A'\)~; il vient alors :
\[
 (n,1+X_1+\ldots+X_n)
 =
 \left(n,\prod_{i=1}^n(1+T_i)\right)
 =
 \sum_{i=1}^n(1,1+T_i)
\]
On en déduit, d'après (6.2.1) :
\[
  \lambda_t(x)=\prod_{i=1}^n\bigl(1+(1,1+T_i)t\bigr),
\]
et on en tire :
\[
  \lambda^k(x)=0\ \text{pour }k>n,\quad\text{et}\quad
  \lambda^n(x)=\prod_{i=1}^n(1,1+T_i)
   =\left(1,1+\sum_{i=1}^nT_i\right)=(1,1+X_1)
\]
Enfin,
\[
  \sum_{i=0}^n(-1)^i\lambda^i(x)=\lambda_{-1}(x)
  =
  \prod_{i=1}^n\bigl(1-(1,1+T_i)\bigr)
  =
  (-1)^n.\prod_{i=1}^n(0,1+T_i).
\]
Or, d'après (6.5.3), \(\prod_{i=1}^n(0,1+T_i)=\left(0,1+(-1)^{n-1}(n-1)!\prod_{i=1}^nT_i+\ldots\right)\)
soit : \(\lambda_{-1}(x)=(0,1-(n-1)!X_n+\ldots)\).

\textbf{Corollaire 6.6.2.} \uline{La relation} 6.6.1.iii) \uline{est vraie pour tout élément} \(X\) \uline{d'augmentation} \(n\).

Pour le montrer nous utiliserons le lemme :

\textbf{Lemme 6.6.3.} \uline{Pour tout} \(n\), \((\Ch A)_n\) \uline{est un} \(\lambda\)-\uline{idéal de} \(\Ch A\).

On peut supposer \(n\ge1\). Soit \(x=(0,\varphi)\in(\Ch A)_n\)~; on peut écrire \(\gamma^i(0,\varphi)=(0,1+Q''_{i,1}+\ldots+Q''_{i,j}+\ldots)\) les \(Q''_{i,j}\) étant les composantes isobares de poids \(j\) des \(Q'_{i,j}\) (3.6). Pour tout \(j\), les polynômes \(Q''_{i,j}\) sont donc des polynômes par rapport aux \(j\) premiers coefficients de \(\varphi\). Si \(j<n\), ils sont donc nuls, et par suite, les \(\gamma^i(x)\) appartiennent à \((\Ch A)_n\), et 
